\SourceLeaf{198}
\begin{center}
{\large ALLGEMEINE AUFLÖSUNG DER AUFGABE}\par\bigskip
DIE THEILE EINER GEGEBENEN FLÄCHE\par\medskip
AUF EINER ANDERN GEGEBENEN FLÄCHE SO ABZUBILDEN\par\medskip
DASS DIE ABBILDUNG DEM ABGEBILDETEN\par\medskip
IN DEN KLEINSTEN THEILEN ÄHNLICH WIRD.\par\bigskip
\rule{.25\linewidth}{.3pt}
\end{center}
\renewcommand{\Article}[1]{\par\medskip\Needspace{4\baselineskip}\pdfbookmark[1]{Art. #1}{abbildung-art-#1}\centerline{\textbf{#1.}}\nobreak\smallskip}
\Article{1}
Die Natur einer krummen Fläche wird durch eine Gleichung zwischen den sich auf jeden Punkt derselben beziehenden Coordinaten \(x,y,z\) bestimmt. Vermöge dieser Gleichung kann jede dieser drei veränderlichen Grössen wie eine Function der beiden andern betrachtet werden. Noch allgemeiner ist es, noch zwei neue veränderliche Grössen \(t,u\) einzuführen, und jede der \(x,y,z\) als eine Function von \(t\) und \(u\) darzustellen, wodurch, wenigstens allgemein zu reden, bestimmte Werthe von \(t\) und \(u\) allemal einem bestimmten Punkte der Oberfläche angehören, und umgekehrt.

\Article{2}
In Beziehung auf eine zweite krumme Fläche sollen \(X,Y,Z,T,U\) ähnliche Bedeutungen haben, wie resp. \(x,y,z,t,u\) in Beziehung auf die erstere.

\Article{3}
Die erste Fläche auf der zweiten \textit{abbilden} heisst, ein Gesetz festsetzen, nach welchem einem jeden Punkte der ersten Fläche ein bestimmter Punkt der zweiten entsprechen soll. Dieses wird dadurch geschehen, dass \(T\) und \(U\) bestimmten Functionen der zwei veränderlichen Grössen \(t\) und \(u\) gleich gesetzt werden.
